\emph{By Brian Champ and Michelle Robidoux}

As 2020 began, a struggle by Indigenous Wet'suwet'en people against attempts to push a natural gas pipeline through their unceded territory sparked a spectacular wave of solidarity across Canada.\footnote{Unceded territory refers to lands that Indigenous people never ceded or legally signed away to the Crown or to Canada.} For almost a month, freight and passenger railway traffic ground to a halt as the \#ShutDownCanada movement spread across the country.\footnote{This article is dedicated to the memory of Philip Murton, a socialist and trade-union activist who died as this article was being written. Thanks to Merv King, member of the Bear Clan of Timiskaming First Nation and activist in the United Steelworkers, and many others for comments on the original draft. All errors and omissions are ours.} The dynamic of Indigenous resistance to encroachments on their sovereignty and mass solidarity by non-Indigenous people---including organised workers---captured the imagination of thousands looking for a strategy in the face of an accelerating climate crisis.

This movement takes place in a global context of struggles confronting the destructive practices of extractive industries. As Martin Upchurch has noted in a previous edition of this journal, ``Many protests have succeeded against the power of the oil companies because of the development of militant tactics and mass support\ldots The number of such protests appears to be growing on a global scale, with more than 3,100 so far reported on the `Global Atlas of Environmental Justice\,'''.\footnote{Upchurch, 2020.} These struggles target key ``chokepoints'' of the capitalist economy, blockading pipeline construction, highways, railways and ports. Indigenous people are at the forefront of these struggles, from Australia to South America and Africa. On Turtle Island---the name given by some Indigenous peoples to the North American continent---this is visible in the battle at Standing Rock against the Dakota Access Pipeline, the successful fight against Enbridge's Northern Gateway Pipelines, the mass opposition to the Energy East pipeline, and the victory of the Inuit people of Clyde River against seismic testing for oil in the Arctic.

These struggles have transformed the landscape, raising hopes of stopping rapacious megaprojects and the climate chaos they bring. They also outline a path along which a powerful movement uniting Indigenous peoples and non-Indigenous settlers in countries like Canada could strike at the roots of capitalism itself.

Yet with every advance in stopping such projects, right-wing forces attempt to widen divisions between Indigenous peoples, climate activists and the broader multi-racial working class. Those divisions, underpinned by centuries of colonialism that has structured the relationships between Indigenous people and settlers, can suddenly explode. In mid-September, when Mi'kmaq fishers exercised their treaty rights by placing lobster traps in the Bay of Fundy, they were set upon by non-Indigenous fishers who rammed their boats, cut their traps, destroyed their vehicles and set fire to a lobster pound. Fuelling these attacks was the claim that Indigenous fishers threaten the sustainability of lobster stocks.\footnote{Edwards, 2020.} This false claim has been allowed to stand by the federal government because it suits them to divide fishers. In reality, although the ten Sipekne'katik lobster licences authorise only 50 traps per licence, a single commercial harvester can deploy up to 400 traps.\footnote{APTN National News, 2020.} The federal government has sat back and let this rendition of colonial violence play out, as police watch it unfold. Tragically, most of the non-Indigenous fishers carrying out these attacks are French-speaking Acadians, whose ancestors only survived hunger, disease, a harsh climate and the British scorched-earth policy because of assistance from the Mi'kmaq.

The Wet'suwet'en struggle and that of the Mi'kmaq community of Sipekne'katik reflect the polarisation that accelerating climate chaos, economic crisis and the Covid-19 pandemic have intensified. The ongoing attacks on Indigenous peoples are a daily reminder that the colonial violence and cruelty that Canada was founded upon continues today. Any movement that hopes to transform society must confront the racism baked into the entire structure of capitalist Canada. The ruling class's deployment of age-old colonial tactics reflects the reality of their limited options as these crises bear down. Even as millions around the world are fighting for a Green New Deal and a ``just recovery'' from the Covid-19 crisis, the Canadian state has doubled down on tar-sands oil. This can only mean confrontation---with Indigenous peoples whose sovereignty is an obstacle to hydrocarbon developments, and with the vast majority of Canadians who see a move away from fossil fuels as necessary.\footnote{Abacus Data, 2020.} Anti-Indigenous racism is used to justify ongoing Canadian colonialism, scapegoating Indigenous people for the horrific conditions they face. Understanding the dynamics at play, their historic roots and the means by which Indigenous resistance and workers' struggles can link up is crucial for the success of both. This article hopes to contribute to that understanding.

\paragraph{Settler colonialism on Turtle Island}

The settler-colonial states of Canada, the United States, Australia and New Zealand, as they exist today, share a majority Anglo-European culture because colonialism decimated the Indigenous populations that have existed for thousands of years. Even as new waves of immigration change the make up of the Canadian population, the result of historic and ongoing genocidal policies mean that Indigenous peoples are now a small minority on their ancestral lands. As Roxanne Dunbar-Ortiz writes:

\blockquote{Nearly all the population areas of the Americas were reduced by 90 percent following the onset of colonising projects, decreasing the targeted Indigenous populations of the Americas from one hundred million to ten million. Commonly referred to as the most extreme demographic disaster---framed as natural---in human history, it was rarely called genocide until the rise of Indigenous movements in the mid-20th century forged questions.\footnote{Dunbar-Ortiz, 2014, p39.}}

Yet despite facing fierce oppression, Indigenous people have played an enormous role in key environmental battles, defending their land against destructive megaprojects. Through centuries of resistance, Indigenous people have won forms of land tenure and legal standing within the colonial framework. Corporations or governments that want to set up a mine, build a pipeline or engage in logging on Indigenous land can't ignore this standing. On Turtle Island, militant ecologism is an organic part of many Indigenous struggles. This outlook stems partly from a continuing relationship to the land, as many Indigenous peoples still rely on hunting, fishing and other subsistence activities. However, although many still retain strong connections to the land, over half of Indigenous people in Canada live in urban centres. The brutal policies of dispossession that created and maintain the Canadian state have severed many Indigenous people's relationship to the land. Organic community memory of the ideas of egalitarianism and reciprocity that characterised pre-colonial Indigenous societies has been reinforced and reinvigorated through struggle. The assertion of traditions and cultures that stand in stark contrast to capitalism inspires many activists who reject the unchecked destruction produced by capitalism and view decolonisation and solidarity as key to challenging this.\footnote{See Lukacs, 2020.}

\paragraph{``Shut Down Canada''}

The electrifying movement to \#ShutDownCanada arose in response to the invasion of Indigenous lands in the north of the province now called British Columbia by the Royal Canadian Mounted Police (RCMP). The police violently arrested Wet'suwet'en matriarchs and supporters. At issue was Indigenous resistance to Coastal GasLink's construction of a multi-billion-dollar pipeline that would channel fracked gas from the interior of British Columbia to export facilities on the Pacific coast. The Wet'suwet'en people have inhabited their territory for at least 10,000 years, and Wet'suwet'en territory encompasses 22,000 km\textsuperscript{2} of northern British Columbia. For years, energy corporations have sought to build pipelines from Alberta through this area of British Columbia in order to ship their products to lucrative markets. The Wet'suwet'en hereditary chiefs have consistently opposed these projects---from the now-defeated Northern Gateway Pipeline to the Pacific Trails Pipeline and the Coastal GasLink project.

On New Year's Eve, 2019, the British Columbia Supreme Court granted Coastal GasLink an injunction to access Wet'suwet'en land in order to continue pipeline construction. Coastal GasLink had signed agreements for the pipeline with five of six Wet'suwet'en band councils, whose leaders are chosen by democratic elections on individual reserves. However, the Wet'suwet'en hereditary chiefs have never ceded their territory, and band councils, which have a limited role in governing six reserves totalling just 35 km\textsuperscript{2}, have no authority over Wet'suwet'en territory.\footnote{For a brief overview of Wet'suwet'en governance structure, see Unist'ot'en Camp website at https://unistoten.camp/about/governance-structure}

In response to the Coastal GasLink injunction, the hereditary chiefs invoked their own law and evicted Coastal GasLink employees. This was despite the looming threat of police violence; indeed, a year earlier, the RCMP had deployed snipers as they surrounded land defenders while attempting to enforce another injunction.\footnote{Dillon and Parrish, 2019.} The chiefs called for provincial construction permits to be revoked and for nation-to-nation talks with social-democratic British Columbia premier John Horgan, Liberal Party prime minister Justin Trudeau, and the commissioner of the RCMP. The British Columbia government, a coalition of the Labour-style New Democratic Party and the Green Party, had just become the first in North America to pass legislation that implemented the United Nations Declaration on the Rights of Indigenous People (UNDRIP). These rights include the requirement of ``free, prior and informed consent'' from Indigenous peoples before any project affecting their lands or territories is approved. The ink was not even dry on that legislation when the British Columbia government claimed the situation was out of their hands because the RCMP were enforcing a court-issued injunction. A leaked letter later revealed that the British Columbia Solicitor General had called in the RCMP by declaring a provincial emergency.\footnote{British Columbia Civil Liberties Association, 2020.}

For his part, Justin Trudeau, who since his election in 2015 repeated that ``no relationship is more important to Canada than the relationship with Indigenous people'', instead called for the ``rule of law'', a demand from colonial power that echoes through the generations.\footnote{Hyslop, 2020.} As the hereditary chiefs stated:

\blockquote{Anuc `nu'at'en (Wet'suwet'en law) is not a ``belief'' or a ``point of view''. It is a way of sustainably managing our territories and relations with one another and the world around us, and it has worked for millennia to keep our territories intact. Our law is central to our identity. The ongoing criminalisation of our laws by Canada's courts and industrial police is an attempt at genocide, an attempt to extinguish Wet'suwet'en identity itself.\footnote{Rubin, 2020.}}

On 6 February, RCMP began arresting land defenders. Within hours, the British Columbia legislature was occupied and ports, railways and highways were blocked across the country. This shut down Canadian National Railways's eastern network, bringing freight traffic from Halifax to Toronto to a halt. A solidarity blockade on Tyendinaga Mohawk Territory stopped travel between Montreal and Toronto, one of the busiest rail corridors in the country. Similar actions took place on Listuguj First Nation territory in Quebec, Kahnawake Mohawk territory south of Montreal, and near New Hazelton, British Columbia.\footnote{Forester, 2020.}

On 8 February, protesters blockaded the access road to the Deltaport container terminal, staying overnight. This port, located in Greater Vancouver, sees \$1 billion worth of goods transported through it yearly. The 300 members of International Longshore and Warehouse Union Local 502 who arrived for the morning shift treated the blockade as a picket line and refused to cross, shutting down the port.\footnote{Boynton, 2020.} On 10 February, police enforcing an injunction arrested 14 protesters. Another 43 were arrested in order to reopen the Port of Vancouver.\footnote{Jacques, 2020.} Mass solidarity actions were organised in cities across the country; 15,000 joined a march in Toronto, and there were bridge, port and road shutdowns in Vancouver.

Two weeks into the shutdown, the Canadian Manufacturers and Exporters body estimated that \$425 million in goods were being stranded daily. The business lobby spoke in increasingly frenzied terms of ``catastrophe'' for the economy. By 21 February, the Atlantic Container Line was diverting shipments from Halifax to New York and Baltimore. ``In Montreal, some 4,000 containers sit immobilised on the docks and bulk agricultural products such as grain can no longer reach the port''.\footnote{\emph{Borneo Bulletin}, 2020.} In Vancouver, goods waiting to be shipped east led to a backlog of 50 ships unable to unload their cargoes.\footnote{Reynolds, 2020.}

As solidarity actions rippled across the country and internationally, politicians pleaded with---and threatened---Indigenous land defenders and their supporters. Trudeau warned, ``The situation as it currently stands is unacceptable\ldots Canadians have been patient. Our government has been patient, but it has been two weeks and the barricades need to come down now\ldots The injunctions must be obeyed''.\footnote{Connolly, 2020.} The hereditary chiefs responded, ``We heard Prime Minister Trudeau just a little while ago talking about the inconvenience Canada has suffered. However, there is a difference between inconvenience and injustice''.\footnote{Scherer and Ljunggren, 2020.}

As these threats accumulated, statements of solidarity with the Wet'suwet'en hereditary chiefs rolled in from the British Columbia Government Employees Union (BCGEU), the Canadian Union of Public Employees, the British Columbia Teachers' Federation, postal workers and nurses' unions.\footnote{See Unist'ot'en Camp, 2020.} Canadian Union of Public Employees (CUPE) national president Mark Hancock stated, ``We would never accept this kind of RCMP behaviour towards striking workers on a picket line. Protest is a fundamental right, and the Wet'suwet'en people have a right to protect their unceded territory.'' BCGEU advised that ``members who encounter a picket line at their worksite are expected to exercise their collective agreement and Charter rights to respect that picket line''.\footnote{British Columbia Government Employees Union, 2020. Legal rights are codified in Canada's Charter of Rights and Freedoms.} Ontario CUPE wrote:

\blockquote{While the epicenter of this struggle takes place in this country's westernmost province, we are witnessing and partaking in solidarity actions across the nation. Here in Ontario, allies are blocking rail lines, collecting funds and organising local events to raise awareness. Respecting Indigenous sovereignty and the fight for environmental justice is fundamental for all Canadians.\footnote{CUPE Ontario, 2020.}}

The British Columbia Teachers' Federation stated on behalf of its 45,000 union members:

\blockquote{We stand in solidarity with the Wet'suwet'en peoples and demand that the governments of British Columbia and Canada uphold their responsibilities as laid out in the Supreme Court's Delgamuukw-Gisday'wa decision of 1997. We stand as witnesses at this historic moment when our governments must make a choice to uphold this court decision or continue the ongoing legacy of colonisation.\footnote{British Columbia Teachers' Federation, 2020.}}

This points to a key contradiction confronting the Canadian settler-colonial state. In its landmark decision, the Supreme Court recognised the Wet'suwet'en hereditary chiefs as the rightful decision makers on their traditional territories, or ``yin'tah''. The case of Delgamuukw versus the Queen found that aboriginal title in British Columbia was never extinguished. As Christopher Roth explains:

\blockquote{British Columbia is a land, for the most part, without treaties\ldots This lack of indisputable cession has been significant in Canadian law because of the Royal Proclamation of 1763, in which King George III declared that title to Indian territory was not to be considered extinguished or transferred merely by conquest or occupation but only through voluntary cession\ldots Under this dispensation, British Columbia is the one area of what was called British North America where treaties were utterly neglected as an instrument of colonisation.\footnote{Roth, 2002, p150. The Royal Proclamation was itself in response to resistance by the Indigenous confederacy led by Odawa chief Obwandiyag (Pontiac) that briefly seized ten British military posts in the Great Lakes area in the spring of 1763.}}

The Delgamuukw ruling, he continues, ``amounts to a finding that what we know as British Columbia has not been ceded to the Crown and is thus not part of Canada\ldots But if so-called British Columbia is not Canada, then what kind of territory is it? What lies `underneath,' to borrow the geological metaphors of the legal jargon, can only be aboriginal title''.\footnote{Roth, 2002, p144.}

The 1997 Delgamuukw decision provoked panic among corporations and their political facilitators:

\blockquote{According to one memo, detailing a meeting that took place one day after the Delgamuukw ruling, the then vice-president of the British Columbia Council of Forest Industries, Marlie Beets, remarked that she had spent the previous hour ``trying to calm'' the chief executives she represented. ``Delgamuukw has only created more uncertainty and we are very concerned by how governments will react to the court's findings,'' Beets said. ``The decision makes the need for certainty through surrender all the more clear. We see no other alternative''.\footnote{Lukacs and Pasternak, 2020.}}

This has been the continuous response to legal recognition of aboriginal title by Canada's courts: corporate demands that politicians obtain Indigenous peoples' ``surrender'' and cession to counter ``investment uncertainty''. The Coastal GasLink injunction was granted as if the constitutional question of Wet'suwet'en aboriginal title did not exist. This is a pattern repeated across the country. Indigenous peoples' sovereignty over their lands is trampled, and industries rely on state violence to ensure access to resources.

As Karla Tait, a Unist'ot'en activist arrested on 10 February, stated:

\blockquote{The violence of Canada and the emptiness of its commitments to us as Indigenous people were laid bare in their actions, and their forcible removal of us in the midst of a ceremony\ldots The events we experienced leading up to February 10 are another culmination of the efforts of Canada to discredit and criminalise us for simply existing on our land as we always have---for defending our rights to our landbase and our historic economy, our ways of empowering ourselves---because of the threat to the capitalist economy.\footnote{Brown and Bracken, 2020.}}

\paragraph{Canada's colonial past and present}

This refusal to recognise Indigenous sovereignty is embedded in Canada's foundations, beginning with the ``doctrine of discovery''. As Jennifer Reid explains:

\blockquote{The doctrine of discovery was the legal means by which Europeans claimed rights of sovereignty, property and trade in regions they had allegedly discovered during the age of expansion. These claims were made without consultation with the resident populations in these territories---the people to whom, by any sensible account, the land actually belonged. The doctrine of discovery is a critical component of historical relations between Europeans, their descendants and Indigenous peoples\ldots It is not simply an artifact of colonial history. It is the legal force that defines the limits of all land claims to this day and, more fundamentally, the necessity of land claims at all. To call it into question, even now, would change the rules of the argument entirely. As one journalist put it, ``It is the federal and provincial governments of Canada who are trying to make a claim to land, a claim based on the doctrine of discovery''.\footnote{Reid, 2010.}}

A vast array of measures was deployed to enforce this framework in the interests of Canadian capitalism as it expanded westward through the Prairies region and into what is now British Columbia. By the mid-1800s, ``US expansion had begun to alarm the leaders of British North America.'' The US and Canadian colonists had already squabbled over territory in the War of 1812, but ``now the Americans were pushing west, threatening to swallow up the bulk of North America. The colonial elite in London and Canada became increasingly obsessed with the idea of connecting the Canadian colonies to British Columbia and asserting their claim to all of the territory in between''.\footnote{Shipley, 2020, p33.}

The development of a transcontinental Canada depended on breaking Indigenous peoples' connection to the land: ``As long as the Plains Indians remained strong, and capable of defending their rights and their lands, they presented, together with the M\'{e}tis, a formidable opposition to the imperialist dreams of Ottawa and London''.\footnote{Adams, 1989, p60. The M\'{e}tis are a distinct Indigenous people in Canada and parts of the US, originally the offspring of Indigenous women and European fur traders.} John A Macdonald, Canada's first prime minister, was also the longest serving Minister of Indian Affairs and the architect of Canada's ``Indian'' policy, embodied in the 1876 Indian Act. As he wrote, ``Indian matters and the land granting system form so great a portion of the general policy of the government that I think it necessary for the prime minister, whoever he may be, to have that in his own hands''.\footnote{Hopper, 2018.} He only gave up this portfolio in 1888, once the Canadian Pacific Railway was built, mass settlement of the West was underway, and Indigenous peoples had been decimated, their nations dismembered and confined to tiny allotments of land known as reserves. This was accomplished through horrific violence, starvation, unscrupulous treaty-making, the creation of the cruel residential schools network, the theft of Indigenous children, the imposition of a tyrannical pass system, and the destruction of Indigenous governance structures and spiritual practices. These methods, which Dene scholar Glen Coulthard describes as a ``genocidal exclusion/assimilation double'', sought to destroy Indigenous communities as distinct peoples.\footnote{Coulthard, 2014, p6. The Dene are an Indigenous group in Canada's arctic and boreal regions.}

Resistance to this assault was continuous and fierce. To enforce its policies, the newly founded Canadian state established the North-West Mounted Police, a permanent occupation force modelled on the Royal Irish Constabulary.\footnote{Tellingly, ``these colonial troops actively compared Indigenous people to the Irish, one declaring to the raucous laughter of his mates that the Sioux `sound exactly like the Irish'. He followed it up with a mockery of the two accents mashed-up. Ireland, of course, was one of Britain's earliest colonial projects, and the Irish were engaged in a very serious struggle against British colonialism, which was spilling over into Canada.''---Shipley, 2020, p53.} The forerunner of today's RCMP, it was instrumental in suppressing ceremonies such as the Thirst Dance in the Prairies and the Potlatch on the Pacific coast. These ceremonies played an important role in redistributing surplus among members of Indigenous societies. Government officials ``objected to the fact that they undermined the accumulation of private property, took people away from agricultural pursuits, brought together bands that they were trying to keep separate, and strengthened the status of traditional leaders and elders''.\footnote{Truth and Reconciliation Commission of Canada, 2015, p130.}

As Indian commissioner Hayter Reed wrote in 1889, ``the policy of destroying the tribal or communist system is assailed in every possible way.'' This included the creation of residential schools where, as Bishop Vital Grandin said in 1875, ``We instil in them a pronounced distaste for the Native life so that they will be humiliated when reminded of their origin. When they graduate from our institutions, the children have lost everything Native except their blood.'' Over the course of a century, an estimated 150,000 First Nations, Inuit and M\'{e}tis children passed through the residential school system. The 2015 report of the Truth and Reconciliation Commission (TRC) reads like an encyclopedia of sadism. Established in 2008 to record the history and impacts of the residential schools, the TRC documented the deaths of over 6,000 students. The report concluded that the brutality deployed to suppress Indigenous language and culture constitutes ``the darkest, longest, and most chilling chapter in the history of the colonisation of Aboriginal peoples''.\footnote{Truth and Reconciliation Commission of Canada, 2015, p131. Far from being limited to residential schools, this policy was in force through the ``Sixties Scoop'', in which Indigenous children were removed from their families and placed in foster care, and continues in child welfare policies today.}

As Tyler Shipley notes:

\blockquote{The development of Canada fits the predictable patterns of a settler colony born out of European expansion between the 15th to 19th centuries. Though it gradually took on a uniquely Canadian character, many of the same basic patterns are found in studies of the US or Australia, or in more modern cases that follow similar logics, from South African apartheid to Israel's occupation of Palestine.\footnote{Shipley, 2020, p91.}}

Considering the catastrophe that befell Indigenous peoples on Turtle Island, the degree to which their communities have survived and are today resurgent is a testament to their centuries of resistance. Yet this catastrophe is not something of the past. It is embodied in state policies deployed today. As Coulthard explains:

\blockquote{Settler-colonial formations are territorially acquisitive in perpetuity\ldots In the specific context of Canadian settler colonialism though, the means by which the colonial state has sought to eliminate Indigenous peoples in order to gain access to our lands and resources have varied over the last two centuries. These range from violent dispossession to the legislative elimination of First Nations legal status under the sexist and racist provisions of the Indian Act, as well as the ``negotiation'' of what are still essentially land surrenders under the present comprehensive land claims policy. The ends have always remained the same: to shore up continued access to Indigenous peoples' territories for the purposes of state formation, settlement, and capitalist development.\footnote{Coulthard, 2014, chapter 4.}}

\paragraph{The consequences of colonialism}

Indigenous peoples in what is today called Canada constitute 4.5 percent of the overall population and are linguistically, culturally and geographically diverse. The Canadian state's attempt to eradicate Indigenous people is not merely a historic phenomenon---it is a colonial present that manifests in myriad ways. This was captured by the government's 2018 Public Health Agency report, \emph{Key Health Inequalities in Canada---A National Portrait}.\footnote{Public Health Agency of Canada, 2018.} Key findings of the report include:

\begin{itemize}
\item Life expectancy for people from Inuit, First Nations and M\'{e}tis communities is 69.7, 70.5 and 74.8 years respectively, compared to the national average of 81.8 years;
\item Infant mortality rates are 2 to 4 times the national average;
\item The suicide rate in Inuit communities is 6.5 times higher than the national average---a testament to the hopelessness that many Indigenous people feel;
\item Rates of arthritis, asthma, diabetes and disability are significantly higher than the national average, as are rates of infectious disease;
\item First Nations reserves are chronically underserved by health services; and,
\item Indigenous peoples in Canada experience the highest levels of poverty. Some 25 percent live in poverty, including a staggering 40 percent of Indigenous children.
\end{itemize}

In 2019, the UN special rapporteur on adequate housing found that:

\blockquote{In Canada, close to half of all First Nations people live on reserves, and more than 25 per cent of them live in overcrowded conditions, constituting approximately seven times the proportion of non-Indigenous people nationally. More than 10,000 on-reserve homes in Canada are without indoor plumbing, and 25 percent of reserves in Canada have substandard water or sewage systems. In a country with more fresh water than anywhere else in the world, 75 per cent of the reserves in Canada have contaminated water.\footnote{United Nations General Assembly, 2019.}}

As of September 2020, 63 reserves have a long-term public health directive that drinking water should be boiled. Lack of clean water, chronic shortage of decent housing and overcrowding have dire consequences. In a pandemic, they are deadly.

This mass of statistics does not begin to tell the story of the deeply rooted systemic racism towards Indigenous peoples in every aspect of Canadian society. The RCMP, founded to suppress Indigenous resistance on the Plains, continue their brutal and often murderous persecution of Indigenous people. In Saskatchewan, 62.5 percent of people who have died from police encounters were Indigenous, despite being only 11 percent of the population.\footnote{Palmater, 2020, p76.} From 2017 to June 2020, almost 40 percent of people killed by police were Indigenous. In 2016, Indigenous people represented 25 percent of the national male prison population and 35 percent of the female prison population.\footnote{Flanagan, 2020.} But it goes beyond police and prisons. On 28 September 2020, Joyce Echaquan, a 37-year-old Atikamekw woman, broadcast her dying moments from her hospital bed in Joliette, Quebec as she suffered racist abuse by hospital staff. In July 2017, Barbara Kentner, a 34-year-old Anishinaabe woman, died from injuries sustained after she was hit by a trailer hitch thrown from a passing truck. She heard her assailant yell ``I got one!'' as the truck drove away. This racism, built into the foundations of Canada, is systematically promoted from the top of society.

\paragraph{From Idle No More to \#ShutDownCanada}

In November 2012, resistance to legislation under Conservative Party prime minister Stephen Harper sparked a resurgence of Indigenous activism. The legislation threatened Indigenous treaty rights by dismantling environmental protection laws. Idle No More erupted on the political scene, with flash-mob round dances occupying malls and government buildings, city centres and highways. It emerged in the context of a global wave of mass movements challenging the status quo. Janaya Khan, a founder of Black Lives Matter Toronto, writes that ``Idle No More played a pivotal role in normalising the language of `Indigenous sovereignty' across leftist and movement cultures, academia and non-Indigenous communities''.\footnote{Khan, 2020, p123.} The movement emerged in the context of growing calls for an inquiry into missing and murdered Indigenous women; between 1980 and 2012, Indigenous women and girls represented 16 percent of all female homicides, despite making up only 4 percent of the female population in Canada.\footnote{Shipley, 2020, p84.} In Saskatchewan, 55 percent of all murdered women are Indigenous. As the Mi'kmaq lawyer and activist Pam Palmater notes, ``The insidious ways in which racism and sexism combine to target Indigenous women and girls for disproportionate levels of abuse, neglect, exploitation, sexualised violence, disappearances and murders is a national crisis that requires a priority emergency response''.\footnote{Palmater, 2020, p272.} Idle No More also emerged as the Truth and Reconciliation Commission conducted hearings across the country. It shaped the terrain on which the findings of the TRC were received in June 2015, when a poll by Angus Reid Institute found that 7 out of 10 agreed that the treatment of Indigenous children at residential schools amounted to ``cultural genocide''.\footnote{Angus Reid Institute, 2015.}

Harper was defeated in the October 2015 federal election and his successor, Trudeau, centred his ``brand'' on the promise of a new nation-to-nation relationship with Indigenous peoples, as well as a commitment to tackling climate change. After years of Harper, there was a palpable sense of relief that the climate change-denying Conservatives were gone, and that the days of overt assimilationist policies were over. Yet as journalist Martin Lukacs writes:

\blockquote{The transformation underway among the Liberal Party, government institutions and the broader establishment was less a sea change than a shape shift. Faced with an Indigenous uprising unlike anything in Canadian history, they were prepared to accept, and even help construct, a new public consensus---making a taboo of overt racism, cleansing our public squares of ugly tokens of our past, embracing the resurgence in Indigenous cultural expression, and adopting the language of Indigenous liberation. But within this consensus, there were several great unmentionables: land, resources and power, and the sharing of any of them.\footnote{Lukacs, 2019, p137.}}

The contradiction between Trudeau's brand and the reality of Canadian capitalism's commitment to extractive megaprojects was soon impossible to hide. At a 2017 energy conference in Houston, Trudeau declared:

\blockquote{No country would find 173 billion barrels of oil in the ground and just leave them there. The resource will be developed. Our job is to ensure that this is done responsibly, safely, and sustainably.\footnote{\emph{Maclean's}, 2017.}}

In June 2019, the Liberal government passed a non-binding motion declaring that Canada is in a climate emergency. The next day, it approved the Trans Mountain tar-sands pipeline, which it had purchased the previous year for \$4.5 billion CAD (ballooning to \$16 billion by February 2020) and which, if completed, would carry 890,000 barrels of oil a day from Alberta to British Columbia. ``Trudeau's pipeline'' runs through unceded Indigenous territories. First Nations leaders have said it will not be built.\footnote{Beaumont, 2019.} His stated commitment to establishing a ``new relationship'' with Indigenous peoples, and meeting Canada's emissions target under the Paris Climate Agreement, is on a collision course with his material support for extracting and marketing tar-sands oil.

\paragraph{Paths of resistance}

The struggle of Wet'suwet'en people struck a powerful chord of solidarity that succeeded in shutting down Canada---destabilising the ``conditions of certainty'' demanded by investors. This solidarity did not come out of the blue. It has been built over years, and a series of initiatives seeking to create a common framework of struggles around climate issues, Indigenous sovereignty and workers' rights moved it forward. \emph{The Leap Manifesto}, created by Indigenous, climate, labour, faith and social justice organisations, was launched in September 2015, attempting to pressure the reformist New Democratic Party to adopt a bold platform for the upcoming elections. It stated:

\blockquote{We start from the premise that Canada is facing the deepest crisis in recent memory. The Truth and Reconciliation Commission has acknowledged shocking details about the violence of Canada's near past. Deepening poverty and inequality are a scar on the country's present. And Canada's record on climate change is a crime against humanity's future\ldots Climate scientists have told us that this is the decade to take decisive action to prevent catastrophic global warming. That means small steps will no longer get us where we need to go. So we need to leap. This leap must begin by respecting the inherent rights and title of the original caretakers of this land. Indigenous communities have been at the forefront of protecting rivers, coasts, forests and lands from out-of-control industrial activity. We can bolster this role, and reset our relationship, by fully implementing the United Nations Declaration on the Rights of Indigenous Peoples.\footnote{See https://leapmanifesto.org/en/the-leap-manifesto. For a discussion of the global movement of Indigenous peoples that led to the UN Declaration, see Dunbar-Ortiz, 2016, p84.}}

In the summer of 2019, the Pact for a Green New Deal was launched, building on the \emph{Leap Manifesto}. It declared that cutting Canada's emissions in half in 11 years, and ensuring no one is left behind in this transition, ``means beginning with the foundational rights and sovereignty of Indigenous communities and implementing the UN Declaration.'' On 27 September, the global climate strike drew 500,000 people on the streets of Montreal, and 100,000 in Vancouver. Toronto saw 50,000 people turn out, including students, Indigenous activists, union members and many more. They were joined by a contingent of auto workers who were worried about the impending closure of the General Motors plant in Oshawa, Ontario. These workers had developed the Green Jobs Oshawa (GJO) campaign, calling for nationalisation of the plant and its conversion into a factory for electric vehicles. A representative of GJO subsequently spoke at the November climate strike in Toronto, where protests targeted the financial backers of the Coastal GasLink pipeline. This built deeper connections between climate, social justice and Indigenous activists, preparing the ground for the direct action of the new year. The sense of possibility was building, as climate actions took place around the world. Yet in spite of this broad support for a Green New Deal, it proved too ``radical'' for the New Democratic Party. Further protests in January 2020 against the impending invasion of Wet'suwet'en territory prepared the ground for the blockades of February.

In early March, the last railway blockades came down. Then the coronavirus pandemic hit. The railway was shut down again, but this time by Covid-19. As governments began rolling out economic support measures, the \#JustRecoveryforAll initiative was launched. Over 150 organisations signed on, demanding that recovery efforts support a just transition, strengthen the social safety net, ``uphold Indigenous rights, and include the full and effective participation of Indigenous peoples, in line with the standard of free, prior, and informed consent''.\footnote{See https://justrecoveryforall.ca}

Canada's history of lies and denial of the impact of colonisation on Indigenous people has been exploded by Indigenous resistance, which has elicited solidarity from millions of non-Indigenous people. The scale of support for the Wet'suwet'en struggle shows the potential for defending Indigenous sovereignty and challenging hydrocarbon megaprojects---but also, to challenge capitalism itself. As Glen Coulthard writes, ``For Indigenous nations to live, capitalism must die''.\footnote{Coulthard, 2014, conclusion.} This echoes M\'{e}tis activist and scholar Howard Adams, who wrote in \emph{Prison of Grass}:

\blockquote{It is quite impossible for Indians and M\'{e}tis in their present condition to experience any real liberation within the present capitalist society in terms of gaining control of reserves and communities, being free from discrimination and racism, obtaining full employment, acquiring a decent standard of living, or becoming free and equal people. We need to liberate ourselves from the courts, ballot boxes, school system, churches, and all other agencies that command us to stay in ``our colonised place''. This oppression of the Native people is so deeply rooted in the capitalist system that it cannot be completely eliminated without eliminating capitalism itself.\footnote{Adams, 1989, p176.}}

Anti-Indigenous racism permeates the fabric of Canadian society and is enmeshed in the operation of the settler state as a whole. Even if the Canadian state were to settle land claims, halt pipeline projects and shift the entire energy grid to renewables, this racism would not end. The very existence of Indigenous peoples exposes Canada's dirty secret: that it is built on stolen land. The ongoing poverty and shocking conditions of life faced by Indigenous people in one of the richest countries in the world is a daily rebuke to the Liberal Party's pretences of ``reconciliation''.

Striking at the roots of Indigenous oppression requires overturning capitalism. This would simultaneously challenge the basis of environmental destruction, poverty and exploitation. Amid the multiple crises gripping the world, and the resistance summoned by them, this is a concrete question. The struggles today can begin to lay the basis for a movement powerful enough to tear up the whole system.

A revolutionary strategy looks to the power of Indigenous mobilisations against colonialism and racism allied with the power of the broader working class to shut down the system. United resistance by Indigenous peoples across Canada highlights the shared conditions created by the colonial-capitalist state and provides concrete opportunities for working-class solidarity with Indigenous people. It has brought the struggle into the consciousness of millions and has shown that the racism produced by the system can be challenged by Indigenous and non-Indigenous together.

Zach Wells, a shop steward on the VIA passenger rail network, publicly spoke out in support of the Wet'suwet'en struggle despite the layoffs it caused in his industry:

\blockquote{I wanted to make clear that solidarity can't begin and end with our co-workers. The labour movement is weakened enormously if it doesn't extend solidarity to broader social justice issues. We have more common cause with First Nations than we do with the Liberal Party of Canada, Canadian National Rail and Coastal GasLink.\footnote{Robidoux, 2020.}}

The longshore workers at Deltaport were putting this view into practice when they refused to cross the picket line and shut down the port---inflicting real economic damage on the bosses. Trade unionists also joined blockades of rail lines across the country. These are glimpses of what is possible. Workers need to be won to defending Indigenous peoples' rights if they are to combat the legacy of colonial racism, and Indigenous people need to be won to the power of the working class in order to overturn capitalism.

A strategy that articulates this approach is needed to counter the efforts by the ruling class to weaponise racism against Indigenous people. This necessarily means working through the question of national oppression as it applies to Canada, a country founded on dispossession of Indigenous peoples and the conquest of the French colony of New France. The idea that non-Indigenous working class people could be won to political recognition of the equal national rights of Indigenous peoples, let alone to the revolutionary overthrow of capitalism, may sometimes seem implausible because of the grip of racist and colonial attitudes. However, Indigenous lands were not stolen because of racism; rather, racism was required for the theft of Indigenous lands. As Chris Harman wrote:

\blockquote{Racism did not emerge all at once as a fully formed ideology\ldots The early attitude to the native inhabitants of North America tended to be that they differed from Europeans because they faced different conditions of life. Indeed, one problem facing the governors of Jamestown, Virginia was that Indian life had a considerable attraction for white colonists, and ``they prescribed the death penalty for running off to live with Indians''.\footnote{Harman, 1999, p253. For a fuller discussion of the early encounters of settlers and Indigenous peoples in North America, see Empson, 2014, pp58-81.}}

As capitalism developed, many of the people who made up the settler population were driven from their land through the enclosure of commons. As Dunbar-Ortiz explains:

\blockquote{The institutions of colonialism and methods for relocation, deportation and expropriation of land had already been practiced, if not perfected, by the end of the 15th century. The rise of the modern state in Western Europe was based on the accumulation of wealth by means of exploiting human labor and displacing millions of subsistence producers from their lands.\footnote{Dunbar-Ortiz, 2014, p33.}}

She continues:

\blockquote{Denied access to the former commons, rural subsistence farmers and their children had no choice but to work in the new woollen textile factories under miserable conditions---that is, when they could find such work, as unemployment was high. Employed or not, this displaced population was available to serve as settlers in the North American British colonies, many of them as indentured servants, motivated by the promise of land. After serving their terms of indenture, they were free to squat on Indigenous land and become farmers again. In this way, surplus labor created not only low labor costs and great profits for the woolens manufacturers but also a supply of settlers for the colonies, which was used as an ``escape valve'' for the home country, where impoverishment might lead to uprisings of the exploited.\footnote{Dunbar-Ortiz, 2014, p35.}}

Settler-colonial capitalism continues to rely on the concept of ``terra nullius'', the idea that land was unowned before colonisation, and the doctrine of discovery to justify its historic and continued oppression of Indigenous people. Although governments and corporations proclaim their commitments to ``reconciliation'', they simultaneously reproduce racist structures and the ideology that sustains them. However, common struggles can challenge the hold of these ideas.\footnote{On Marxism and the national question, see Karl Marx, \emph{The Irish Question} and Lenin's \emph{The Right of Nations to Self-Determination}. See also Crouch, 2003.} Describing the struggle against Dakota Access Pipeline at Standing Rock, Lakota scholar and activist Nick Estes writes:

\blockquote{Political elites and corporate media have frequently depicted poor whites and poor Natives as irreconcilable enemies, lacking common ground and competing for scarce resources in economically depressed rural areas. Yet the defence of Native land, water and treaties brought us together. Although not perfect, the Oceti Sakowin camp was a home to many for months. And the bonds were long lasting, despite the horrific histories working against them.\footnote{Estes, 2019, p7.}}

\paragraph{Imagining a world beyond capitalism}

The Canadian philosopher James Tully has written:

\blockquote{Practices of reconciliation are doomed to failure if they attempt to change Indigenous-settler relations without dismantling vicious economic relations or if they attempt to change global economic systems without addressing the exercise of colonial power over Indigenous peoples worldwide\ldots Yet this is unfortunately the approach offered in British Columbia and elsewhere in Canada by mainstream models of Indigenous-settler reconciliation, with their ``business as usual'' imperative.\footnote{Tully, 2020, p402.}}

This approach has been embraced by all mainstream political parties in Canada, including the New Democratic Party. The right-wing government of Alberta, the heart of the tar-sands industry, has invested billions in pipeline development despite corporate investment drying up with the collapse of oil prices. This includes \$1.5 billion of public funds invested in the Keystone XL pipeline, whose construction US president Joe Biden has now cancelled. Alberta premier Jason Kenney has introduced laws criminalising blocking ``critical infrastructure'' and used Covid-19 emergency measures to suspend labour and environmental legislation. These investments and laws protect the oil companies, the forces of death on a planetary scale. As well, far-right groups increasingly threaten Indigenous land defenders and pipeline opponents. This adds to the importance of building solidarity among the wider population, including workers. In the midst of the pandemic, the Alberta government launched a major attack on healthcare workers, announcing 11,000 job losses through outsourcing laundry, food services, housekeeping and laboratory jobs. Despite Covid-19 conditions, wildcat strikes swept through 49 worksites in 39 cities and towns on 26 October. This is the type of action needed to resist the unfolding disasters of capitalism.

That is why initiatives like the \emph{Leap Manifesto}, the Green New Deal and \#JustRecoveryforAll are so critically important. They link demands for Indigenous self-determination with those for a just transition for workers. Such bold proposals can provide an alternative to working-class identification with the settler-colonial state, towards solidarity with Indigenous people. Workers building pipelines and working in the tar sands and refineries need a concrete alternative to the dead end of fossil fuel production. Indigenous peoples need solidarity from working-class people to win their struggles. This counters the false dichotomies between jobs, Indigenous sovereignty and the environment.

The deep economic, environmental and social crises confronting humanity raise the question of what Karl Marx and Frederick Engels describe in the \emph{Communist Manifesto} as a fight that ends ``either in a revolutionary reconstitution of society at large, or in the common ruin of the contending classes''. They advance a vision of socialism where the vast majority---the ``associated producers''---wrests control of society away from the minority ruling class in order to democratically and cooperatively organise the economy to meet human need instead of corporate greed.

In \emph{Marx's Ecology}, John Bellamy Foster writes:

\blockquote{The most important problem facing the society of associated producers, Marx emphasised again and again in his work, would be to address the problem of the metabolic relation between human beings and nature under the more advanced industrial conditions prevailing in the wake of the final revolutionary crisis of capitalist society. To this end, it was clearly necessary to learn more about the human relation to nature and subsistence, through the development of property forms, over the great span of ethnological time. Marx was thus driven back, by the materialist precepts of his analysis, to a consideration of the origins of human society and the human relation to nature---as a means for envisioning the potential for a more complete transcendence of an alienated experience.\footnote{Foster, 2000, p221.}}

In the last years of his life Marx was captivated by the study of Indigenous societies.\footnote{See Anderson, 2010; also Foster, Clark and Holleman, 2020.} In his \emph{Ethnological Notebooks}, he carefully pored over writings in the new field of anthropology, including those of Lewis Henry Morgan. As Frank Rosemont writes:

\blockquote{On page after page Marx highlights passages wildly remote from what are usually regarded as the ``standard themes'' of his work. Thus we find him invoking the bell-shaped houses of the coastal tribes of Venezuela; the manufacture of Iroquois belts ``using fine twine made of filaments of elm and basswood bark''; ``the Peruvian legend of Manco C\'{a}pac and Mama Ocllo, children of the sun''; burial customs of the Tuscarora; the Shawnee belief in metempsychosis; ``unwritten literature of myths, legends and traditions''; the ``incipient sciences'' of the village Indians of the South West; the Popul Vuh, sacred book of the ancient Quich\'{e} Maya; the use of porcupine quills in ornamentation; Indian games and ``dancing as a form of worship''.\footnote{Rosemont, 1989, p205.}}

The \emph{Ethnological Notebooks} reveal his deep interest in the concrete experience of societies not distorted by the ravages of capitalism. Indigenous peoples have suffered most acutely from these ravages. However, they have never stopped resisting the system's genocidal logic and through that resistance have held on to ways of life, systems of law, forms of social organisation and cultures that possess a relationship to the natural world that is in sharp contrast to the logic of capitalism. Their continued existence is a testament to longstanding struggles against attempts to assimilate and eradicate them. Importantly, today there is much wider support for these struggles.

Continuing to build links between workers' struggles in the workplace and the ongoing resistance of Indigenous people to the colonial state is crucial for building a path forward for both. They are linked because they are both directed against the capitalist system, suggesting a potential revolutionary convergence that can overthrow capitalism in the context of settler-colonial states. And the battles that gesture towards the need for this type of unity are multiplying.
